\documentclass[11pt]{article}

\usepackage[utf8]{inputenc}
\usepackage[T1]{fontenc}
\usepackage{lmodern}
\usepackage{geometry}
\usepackage{amsmath,amssymb,amsfonts,amsthm,mathtools}
\usepackage{bm}
\usepackage{enumitem}
\usepackage{microtype}
\usepackage{hyperref}
\usepackage{titlesec}
\usepackage{booktabs}
\usepackage{array}
\usepackage{setspace}
\usepackage{mathrsfs}
\usepackage{physics}

\geometry{margin=1in}
\hypersetup{colorlinks=true, linkcolor=black, urlcolor=blue, citecolor=black}
\setlist[enumerate]{topsep=0.4em,itemsep=0.35em}
\setlist[itemize]{topsep=0.3em,itemsep=0.25em}
\titleformat{\section}{\large\bfseries}{\thesection.}{0.5em}{}
\titleformat{\subsection}{\normalsize\bfseries}{\thesubsection.}{0.5em}{}
\setstretch{1.06}

\newcommand{\Aether}{\AE{}ther}
\newcommand{\Aflow}{\AE{}ther-flow}
\newcommand{\TheoryName}{\Aflow{} Interpretation of Relativity}
\newcommand{\TheoryProgram}{\Aether{} / \Aflow{} framework}
\newcommand{\CompletedMetric}{g^{\sharp}}

\newcommand{\CurrentBodyImageCore}{\mathfrak S^{\mathrm{disp,resp,cbimg}}}
\newcommand{\VisibleResponseOperator}{S^{\mathrm{disp,resp,vis}}}
\newcommand{\MinimalReducedStructure}{\mathfrak S^{\mathrm{disp,resp,red,min}}}

\newtheorem{definition}{Definition}
\newtheorem{proposition}{Proposition}
\newtheorem{remark}{Remark}
\newtheorem{corollary}{Corollary}
\newtheorem{theorem}{Theorem}

\title{\textbf{The \TheoryName{}}\\[0.4em]
\large Primitive-Reservoir Observer-Reduction\\
\large Section-Centered Hidden-Response\\
\large Canonical Current-Body Image Core Collapse to\\
\large Canonical Minimal Visible-Response Representative\\
\large Next Benchmark Criterion}
\author{}
\date{}

\begin{document}

\maketitle

\begin{abstract}
The current active-primary theorem now proves that the full canonical current-body image core \(\CurrentBodyImageCore\) already determines a unique positive visible-response operator \(\VisibleResponseOperator\) on the fixed section-centered displacement fibers and therefore collapses benchmark-facingly to the canonical minimal visible-response representative \(\MinimalReducedStructure\). The image body, the restricted projection / canonical-lift relation, and the exact-shell image are no longer independently live benchmark-facing burdens there.

This note records the routing consequence of that gain. The next honest benchmark-facing move must now derive or justify the visible-response operator \(\VisibleResponseOperator\) itself on the fixed displacement fibers, equivalently the canonical minimal visible-response representative \(\MinimalReducedStructure\), from still more primitive explicit substrate structure while preserving the same benchmark one-metric package, the same ordinary matter route, and the same claim boundary. Another manuscript that merely restores the full current image-core package, re-expands to a reduced-carrier presentation, replays the projection-operator core, replays the full displacement-response substrate law, or re-adjoins another same-visible-response carrier-level realization does not count as the clean next step. The one operative metric remains exactly \(\CompletedMetric\), there is no second operative metric, the low-energy matter law is not enlarged, and no stronger promotion language is licensed.
\end{abstract}

\section{Introduction}

The active derivational program of \TheoryName{} is no longer merely at projection-operator-core or reduced-structure scope on the section-centered hidden-response line. The new theorem spends a still smaller option below those frontiers: even the full canonical current-body image core already carries no benchmark-facing content beyond the visible quadratic response on the fixed displacement fibers. What remains live is therefore not the image body, the restricted projection / canonical-lift relation, the exact-shell image, or another same-visible-response realization of those data. It is the visible quadratic response itself.

The present note records the routing consequence of that gain. It does not reopen the full current image-core package, the reduced projection-response substrate structure, the older projection-operator core, the full displacement-response substrate law, or the larger pullback-body-operator, pullback-operator, pullback-quadratic, hidden-response, residual-normal-form, or pre-proto-section presentations as if any of those were again the live burden. It records only which kind of next manuscript would now count as an honest benchmark-facing gain below the current visible-response frontier.

\paragraph{Framework and claim boundary.}
\TheoryProgram{} denotes the broader \Aether{} / \Aflow{} framework built on the claims that the \Aether{} is the underlying four-dimensional substrate of reality, the \Aflow{} is its intrinsic ordered motion, observed three-dimensional space is the local experiential slice of that deeper substrate, S-time is the experienced order of change arising from matter, light, and the \Aflow{}, and observed expansion is the three-dimensional appearance of deeper four-dimensional ordered motion. Within that framework, \TheoryName{} denotes the exact-closure theory in which the effective gravitational dynamics are adopted to be exactly Einsteinian with universal matter coupling. In the active sequence, \emph{adoption} means use of that established relativistic dynamics without claiming substrate derivation, while \emph{derivation} is reserved for a first-principles recovery from explicit substrate variables.

The present note stays strictly inside that boundary. It records only which kind of next manuscript would now count as an honest benchmark-facing gain below the canonical-minimal-visible-response frontier after the current-image-core route has been spent.

\section{Fixed Inputs}

\begin{definition}[Visible-response frontier inputs]
The criterion recorded here uses the following fixed inputs.
\begin{enumerate}
    \item The benchmark exact-closure package, which fixes the public one-metric GR target of the repository.
    \item The benchmark gatekeeping layer, including the recorded safeguard that blocks same-output deeper-origin relays from counting as new front-line progress once the live benchmark-relevant datum is unchanged.
    \item The preserved theorem chain showing that the full canonical current-body image core is forced once a reduced projection-response substrate structure is fixed, that any compatible reduced structure carries the same visible quadratic response operator up to inert kernel extension once the fixed image core is held fixed, and that the older projection-operator-core package already collapses to the same canonical minimal visible-response representative.
    \item The new theorem proving that the full canonical current-body image core \(\CurrentBodyImageCore\) itself already collapses benchmark-facingly to the canonical minimal visible-response representative \(\MinimalReducedStructure\), equivalently to the visible quadratic response operator \(\VisibleResponseOperator\) on the fixed displacement fibers.
\end{enumerate}
\end{definition}

\begin{remark}
The present note does not replace those manuscripts. It records the routing consequence of reading them together under the active safeguard against recursive depth drift.
\end{remark}

\section{Why the Live Burden Now Sits Below Image-Core Restatement}

\begin{proposition}[The live burden is renewed below current-image-core restatement]
At the current stage of the primitive-reservoir observer-reduction line, the current-image-core collapse theorem renews the live benchmark-facing burden below any mere restatement of the same full canonical current-body image core \(\CurrentBodyImageCore\), the same canonical minimal visible-response representative \(\MinimalReducedStructure\), or the same visible quadratic response operator \(\VisibleResponseOperator\) on the fixed displacement fibers.
\end{proposition}

\begin{proof}
That theorem changes benchmark-facing status because it proves that even the full current-image-core package already carries no benchmark-facing content beyond the visible quadratic response on the fixed displacement fibers. Once that collapse is in hand, another manuscript that merely restores the same image body, the same restricted projection / canonical-lift relation, the same exact-shell image, another reduced-carrier realization of the same visible operator, or another older projection/body presentation no longer removes any remaining benchmark-facing freedom. The live burden therefore no longer sits on another realization of the same visible-response datum. It sits on deriving that visible-response datum itself from still more primitive explicit substrate structure.
\end{proof}

\begin{definition}[Benchmark-neutral same-visible-response relay below the current frontier]
A \emph{benchmark-neutral same-visible-response relay below the current frontier} is a manuscript that preserves the same benchmark one-metric output, the same ordinary matter route, the same visible-response operator \(\VisibleResponseOperator\), the same canonical minimal visible-response representative \(\MinimalReducedStructure\), and the same downstream benchmark-facing consequences while merely restoring the already-spent full current-image-core package \(\CurrentBodyImageCore\), re-expanding to a reduced-carrier presentation, re-expanding to the older projection-operator-core or full displacement-response-substrate-law presentation, or otherwise carrying the same visible quadratic response through another larger same-output realization without a new derivation of that visible-response datum itself.
\end{definition}

\section{The Current Post-Image-Core Criterion}

\begin{definition}[Admissible next benchmark-facing manuscript below the current frontier]
Below the current canonical-minimal-visible-response frontier, a manuscript counts as an admissible next benchmark-facing move only if it does at least one of the following:
\begin{enumerate}
    \item derives or justifies the visible-response operator \(\VisibleResponseOperator\) itself on the fixed displacement fibers from still more primitive explicit substrate structure;
    \item equivalently derives or justifies \(\MinimalReducedStructure\) from still more primitive explicit substrate structure;
    \item proves a sharper necessity, uniqueness, or compatibility theorem strictly inside that visible-response datum while preserving the same benchmark one-metric package, the same ordinary matter route, and the same claim boundary.
\end{enumerate}
\end{definition}

\begin{theorem}[Next honest benchmark-facing task below the current frontier]
The next honest benchmark-facing manuscript below the current canonical-minimal-visible-response frontier must derive or justify the visible-response operator \(\VisibleResponseOperator\) itself on the fixed displacement fibers, equivalently the canonical minimal visible-response representative \(\MinimalReducedStructure\), from still more primitive explicit substrate structure, or prove a still sharper theorem strictly inside that visible-response datum. A manuscript that only repeats the same visible-response datum through another restored current-image-core package, another reduced-carrier presentation, another restored projection-operator core, another replay of the full displacement-response substrate law, or another larger same-output presentation does not satisfy the criterion.
\end{theorem}

\begin{proof}
By the preceding proposition, the live burden already lies below any mere restatement of the same visible-response datum or of the same current-image-core realization. Therefore a subsequent manuscript must improve on that visible datum rather than merely accompany it through another representation. A deeper-origin manuscript can do so only if it changes benchmark-facing status by more than presentation alone. The remaining honest alternatives are therefore exactly those stated in the definition.
\end{proof}

\section{What the Next Move Should Not Be}

\begin{proposition}[Screened next moves below the current frontier]
Below the current canonical-minimal-visible-response frontier, none of the following is the default next benchmark-facing move:
\begin{enumerate}
    \item another same-output deeper-origin relay that merely preserves the current visible-response operator \(\VisibleResponseOperator\) and therefore merely reproduces the same canonical minimal visible-response representative \(\MinimalReducedStructure\) together with the same downstream benchmark-facing consequences;
    \item another restoration of the full current-image-core package \(\CurrentBodyImageCore\) as if its image-body, projection / lift, or exact-shell constituents were again independently live;
    \item another restoration of a reduced-carrier presentation, of the older full projection-operator-core body presentation, or of the full displacement-response substrate law as if those already-spent larger packages were again the honest current burden;
    \item another replay of the larger pullback-body-operator, pullback-operator, pullback-quadratic, hidden-response, residual-normal-form, or pre-proto-section presentations as if they were again independently live below the current frontier;
    \item a manuscript that introduces a second operative metric, enlarges the ordinary matter law, or uses stronger promotion language.
\end{enumerate}
\end{proposition}

\begin{proof}
Item (1) fails the preceding criterion because it leaves the renewed live burden unchanged. Item (2) fails because it restores a realization whose benchmark-facing freedom is already screened out by the current theorem. Item (3) fails because it returns to larger already-spent packages rather than deriving the visible-response datum itself. Item (4) likewise re-expands to larger already-spent presentations. Item (5) violates the fixed claim boundary of the benchmark package and therefore cannot count as a disciplined next step on the active line.
\end{proof}

\begin{remark}
This screening statement is not a denial that those deeper or alternative manuscripts may matter strategically. It is a routing statement: they are not the default way to spend the next benchmark-facing move below the current frontier unless they do the same benchmark-facing work named above.
\end{remark}

\section{Conclusion}

The positive result of this note is routing discipline below the canonical-minimal-visible-response frontier after the current-image-core route has been spent. The present theorem remains the live benchmark-facing gain because it already proves that even the full canonical current-body image core carries no benchmark-facing content beyond the visible quadratic response on the fixed displacement fibers. The earlier reduced-structure necessity / uniqueness theorem and the later projection-operator-core collapse theorem remain preserved handoff history because they first showed that compatible reduced-carrier alternatives and then the older projection/body presentation no longer set the live burden either before the present theorem removed the current-image-core realization itself from the honest stopping point.

The scope remains narrow and honest. The benchmark package is unchanged: the one operative metric is still exactly \(\CompletedMetric\), the ordinary matter route is unchanged, there is no second operative metric, and the low-energy matter law is not enlarged. Nothing here upgrades adoption to derivation or licenses stronger promotion language.

The next open burden is therefore clear. The next benchmark-facing manuscript should derive or justify the visible-response operator \(\VisibleResponseOperator\) itself on the fixed displacement fibers, equivalently derive the canonical minimal visible-response representative \(\MinimalReducedStructure\), from still more primitive explicit substrate structure. Another restored current-image-core package, another reduced-carrier replay, another replay of the projection-operator core, or another replay of the full displacement-response substrate law is not the clean next manuscript target at current scope.

\end{document}
